\documentclass[11pt,a4paper]{article}
\usepackage[T1]{fontenc}
\usepackage[utf8]{inputenc}
\usepackage[italian,provide=*]{babel}
\usepackage{mathptmx,microtype}
\usepackage[margin=2.3cm]{geometry}
\usepackage{amsmath,booktabs,graphicx,caption}
\usepackage[hidelinks]{hyperref}
\captionsetup{font=footnotesize,labelfont=bf,skip=3pt}
\graphicspath{{figure/}}
\pagestyle{empty}
\setlength{\parskip}{0.4em}\setlength{\parindent}{0pt}
\begin{document}
\begin{center}
{\Large\bfseries \emph{Volatility risk premium} e strategie non direzionali\\[2pt] in opzioni sull'indice S\&P 500: evidenza dalle \emph{option chain} SPX 2010--2023}\\[6pt]
{\small Novembre 2026}\\[8pt]
{\bfseries Sintesi}
\end{center}
{\let\thefootnote\relax\footnotetext{Autori: Giordano Margotti, Alex Crescentini, Alessandro Celani, Giulio Sciacchitano (gruppo di lavoro \emph{finance-thinking}). Documento di ricerca; non costituisce consulenza in materia di investimenti.}}

\textbf{Domanda.} Il lavoro verifica se nei prezzi delle opzioni europee sull'indice S\&P 500 esista un \emph{volatility risk premium} (VRP) e se le strategie non direzionali di vendita di premio (iron condor, credit spread, strangle e strutture affini) riescano a incassarlo al netto dei costi.

\textbf{Dati e metodo.} \emph{Option chain} \emph{end-of-day} delle opzioni SPX (opzioni su indice di stile europeo, \emph{cash-settled}, quotate al Cboe; fonte OptionsDX) dal gennaio 2010 al dicembre 2023: @@sedute@@ sedute e @@coppie@@ coppie data--scadenza. La volatilità implicita è ricavata dai prezzi con il modello di Black (1976) sul forward implicito nella \emph{put-call parity}; la varianza implicita \emph{model-free} segue la formula del Cboe; la varianza realizzata è misurata fino alla scadenza della stessa opzione. La previsione sotto la misura reale è un GJR-GARCH(1,1) con innovazioni $t$ di Student stimato senza informazioni future. I backtest coprono @@ncfg@@ configurazioni di dieci famiglie, per @@ntrade@@ operazioni simulate, ai prezzi mid e ai prezzi bid-ask.

\textbf{Risultati.} (i) La volatilità \emph{model-free} a 30 giorni supera la realizzata di @@mfvrp30@@ punti in media ($t$ @@mfvrpt30@@); la volatilità \emph{at-the-money} la supera di @@vrp30@@ punti ($t$ @@vrpt30@@). Il premio risiede nelle \emph{wing}, non nel centro della superficie. (ii) Il premio \emph{at-the-money} è significativo nel 2010--2014 (@@pa@@ punti, $t$ @@pat@@) e nullo nel 2015--2019 e nel 2020--2023. (iii) La volatilità implicita \emph{at-the-money} a 30 giorni è una previsione non distorta della realizzata (pendenza @@mzb30@@, test congiunto $p=$ @@mzp30@@) e ha errore inferiore al GARCH e alla volatilità storica. (iv) Le put con delta 0{,}10 a 30 giorni hanno restituito ai compratori una quota pari a @@ins30_10@@ del premio pagato. (v) Lo Sharpe mediano delle configurazioni è @@shmid@@ ai prezzi mid e @@shba@@ ai prezzi bid-ask; la quota con alfa significativamente positivo rispetto all'indice scende dal @@amid@@\% al @@aba@@\%, valore atteso per solo effetto del caso. (vi) I risultati positivi sono concentrati nelle strutture che vendono put e hanno beta azionario positivo; iron condor e straddle hanno risultato netto nullo o negativo; le regole di uscita anticipata riducono il rendimento perché pagano un secondo spread. (vii) La sospensione degli ingressi quando la volatilità realizzata supera il 25\% non migliora lo Sharpe (mediana da @@f0sh@@ a @@f25sh@@) e non riduce il mese peggiore: le perdite estreme hanno origine in fasi di calma, e gli strike scelti per delta si allargano già con la volatilità implicita (dal @@dlo@@\% al @@dhi@@\% del forward tra il terzile basso e il terzile alto). (viii) I rendimenti sono lontani dalla normale: per la short put con delta 0{,}16 a 45 giorni la \emph{skewness} per operazione è @@csp_sk@@ e l'\emph{excess kurtosis} @@csp_ku@@, contro @@ci45_sk@@ e @@ci45_ku@@ per l'indice sullo stesso orizzonte, e la normalità della serie mensile è rifiutata nel @@cq_jb@@\% delle configurazioni. Lo Sharpe sopravvaluta quindi le strutture che vendono volatilità, e il rischio di coda va misurato direttamente.

\begin{center}\small
\begin{tabular}{lrr}\toprule
 & Prezzi mid & Prezzi bid-ask \\ \midrule
Sharpe mediano, @@ncfg@@ configurazioni & @@shmid@@ & @@shba@@ \\
Quota con Sharpe positivo (\%) & @@posmid@@ & @@posba@@ \\
Quota con alfa positivo, $t>2$ (\%) & @@amid@@ & @@aba@@ \\
Short put $\Delta$0{,}16, 45 giorni: Sharpe & @@spsh0@@ & @@spsh1@@ \\
Put credit spread $\Delta$0{,}20/0{,}05, 30 giorni: Sharpe & @@pcssh0@@ & @@pcssh1@@ \\
Iron condor $\Delta$0{,}16/0{,}05, 30 giorni: Sharpe & @@icsh0@@ & @@icsh1@@ \\
\bottomrule\end{tabular}
\end{center}

\textbf{Conclusione.} Il premio esiste come compenso per il rischio di coda ed è coerente con un mercato efficiente sul piano informativo nel centro della superficie. Lo spread bid-ask, l'esposizione azionaria implicita e la selezione a posteriori dei parametri spiegano la quasi totalità dei risultati positivi dei backtest. L'ipotesi che l'incasso sistematico del premio tenda, nel lungo periodo e al netto dei costi, a un risultato nullo non è rifiutata dai dati.

\textbf{Parole chiave:} \emph{volatility risk premium}, opzioni su indice, iron condor, put credit spread, GJR-GARCH, costi di transazione, \emph{multiple testing}.

\newpage
\begin{figure}[t]\centering
\includegraphics[width=0.78\textwidth]{serie_it.pdf}
\caption*{\textbf{Figura 1.} IV \emph{at-the-money} sulla scadenza più vicina a 30 giorni e volatilità realizzata fino a quella scadenza.}
\vspace{3pt}
\includegraphics[width=0.41\textwidth]{struttura_it.pdf}\hfill\includegraphics[width=0.41\textwidth]{hv_it.pdf}
\caption*{\textbf{Figura 2.} A sinistra: medie per orizzonte di IV \emph{model-free}, IV \emph{at-the-money} e volatilità realizzata. A destra: premio \emph{at-the-money} medio successivo per livello della volatilità storica all'ingresso.}
\vspace{3pt}
\includegraphics[width=0.47\textwidth]{sharpe_it.pdf}\hfill\includegraphics[width=0.33\textwidth]{meta_it.pdf}
\caption*{\textbf{Figura 3.} A sinistra: distribuzione dello Sharpe sulle @@ncfg@@ configurazioni ai prezzi mid e bid-ask. A destra: Sharpe nel 2010--2016 e nel 2017--2023 per le configurazioni a credito.}
\vspace{3pt}
\includegraphics[width=0.66\textwidth]{curve_it.pdf}
\caption*{\textbf{Figura 4.} Rendimento cumulato sul margine ai prezzi bid-ask per cinque configurazioni e per l'indice.}
\end{figure}
\end{document}
